\documentclass[11pt]{article}
% ===========================================================================
%  Shared preamble for the AP Calculus AB exam-preparation worksheets.
%  Input right after \documentclass. Files beginning with "_" are not compiled
%  on their own. Navy and gold, to match the course's lesson pages.
% ===========================================================================
\usepackage[letterpaper,top=2.3cm,bottom=2.2cm,left=1.7cm,right=1.7cm,headheight=28pt,footskip=30pt]{geometry}
\usepackage{amsmath,amssymb}
\usepackage[T1]{fontenc}
\usepackage[utf8]{inputenc}
\usepackage{lmodern}
\usepackage{xcolor}
\usepackage[most]{tcolorbox}
\usepackage{enumitem}
\usepackage{fancyhdr}
\usepackage{array}
\usepackage{tikz}
\usepackage{pgfplots}
\pgfplotsset{compat=1.18}

\definecolor{navy}{HTML}{161E54}
\definecolor{ink}{HTML}{27357E}
\definecolor{gold}{HTML}{B98A00}
\definecolor{goldbg}{HTML}{FFFDF3}
\definecolor{inkbg}{HTML}{F6F7FD}
\definecolor{rule}{HTML}{C9CFF0}
\definecolor{okgreen}{HTML}{1B7A44}
\definecolor{okbg}{HTML}{F1F8F3}

\pagestyle{fancy}
\fancyhf{}
\fancyhead[L]{\footnotesize Name: \underline{\hspace{4cm}}}
\fancyhead[C]{\textbf{\color{navy}Integration Academy}\\[-2pt]{\scriptsize Calculus AB \textperiodcentered{} AP\textsuperscript{\textregistered} Exam Preparation}}
\fancyfoot[L]{\tiny\color{black!55}\parbox[t]{0.78\textwidth}{AP\textsuperscript{\textregistered} is a registered trademark of the College Board, which was not involved in the production of, and does not endorse, this material. \textcopyright{} Integration Academy.}}
\fancyfoot[R]{\scriptsize\thepage}
\renewcommand{\headrulewidth}{1.2pt}
\renewcommand{\headrule}{\hbox to\headwidth{\color{gold}\leaders\hrule height \headrulewidth\hfill}}
\renewcommand{\footrulewidth}{0pt}
\setlength{\parindent}{0pt}
\setlength{\parskip}{4pt}

% Title block: \wstitle{code}{title}{unit line}
\newcommand{\wstitle}[3]{\begin{center}{\small\bfseries\color{gold}WORKSHEET #1}\\[3pt]{\LARGE\bfseries\color{navy}#2}\\[3pt]{\small\color{black!60}#3}\end{center}\vspace{2pt}}
% Section heading
\newcommand{\sech}[1]{\par\vspace{10pt}{\large\bfseries\color{navy}#1}\par\nobreak\vspace{2pt}{\color{rule}\hrule height0.8pt}\vspace{7pt}}
% A level tag printed beside a group of questions
\newcommand{\level}[1]{\par\vspace{6pt}{\footnotesize\bfseries\color{gold}\MakeUppercase{#1}}\par\nobreak\vspace{2pt}}

\newtcolorbox{ideabox}{colframe=ink,colback=inkbg,coltitle=white,colbacktitle=ink,fonttitle=\bfseries\small,title={Key ideas},arc=2pt,boxrule=0.8pt,left=8pt,right=8pt}
\newtcolorbox{example}[1]{colframe=rule,colback=white,coltitle=navy,colbacktitle=inkbg,fonttitle=\bfseries\small,title={#1},arc=2pt,boxrule=0.8pt,borderline west={3pt}{0pt}{ink},left=10pt,right=8pt}
\newtcolorbox{frbox}[1]{colframe=gold,colback=goldbg,coltitle=white,colbacktitle=gold,fonttitle=\bfseries\small,title={#1},arc=2pt,boxrule=0.8pt,breakable,left=8pt,right=8pt}
\newtcolorbox{solbox}{colframe=okgreen,colback=okbg,coltitle=white,colbacktitle=okgreen,fonttitle=\bfseries\small,title={Worked solutions},arc=2pt,boxrule=0.8pt,breakable,left=8pt,right=8pt}
\newtcolorbox{scorebox}[1]{colframe=navy,colback=white,coltitle=white,colbacktitle=navy,fonttitle=\bfseries\small,title={#1},arc=2pt,boxrule=0.8pt,breakable,left=8pt,right=8pt}

\newcommand{\soln}{\par\smallskip\textbf{\color{ink}Solution.}\ }
\newcommand{\result}[1]{\par\smallskip\textbf{#1}}
% blank working space of a given height
\newcommand{\work}[1]{\par\vspace{#1}}
% a point value in the margin of a scoring guide
\newcommand{\pts}[1]{\hfill{\footnotesize\bfseries\color{gold}[#1]}}

% Coordinate grid for graph questions: \begin{apgraph}{xmin}{xmax}{ymin}{ymax} ... \end{apgraph}
\newenvironment{apgraph}[4]{\begin{center}\begin{tikzpicture}\begin{axis}[
  axis lines=middle, xmin=#1, xmax=#2, ymin=#3, ymax=#4,
  grid=both, grid style={black!12}, xtick distance=1, ytick distance=1,
  ticklabel style={font=\scriptsize}, xlabel={$x$}, ylabel={$y$},
  width=10cm, height=7.2cm, samples=80, clip=false,
  every axis plot/.append style={line width=1.1pt, color=ink}]}{\end{axis}\end{tikzpicture}\end{center}}
\newcommand{\closeddot}[2]{\addplot[only marks, mark=*, mark size=2.2pt, color=ink] coordinates {(#1,#2)};}
\newcommand{\opendot}[2]{\addplot[only marks, mark=*, mark size=2.2pt, color=ink, mark options={fill=white}] coordinates {(#1,#2)};}

% Multiple-choice options: four across, or stacked when the options are long.
\newcommand{\opts}[4]{\par\vspace{3pt}\begin{tabular}{@{}p{0.235\linewidth}p{0.235\linewidth}p{0.235\linewidth}p{0.235\linewidth}@{}}(A)\ #1 & (B)\ #2 & (C)\ #3 & (D)\ #4\end{tabular}\par\vspace{9pt}}
\newcommand{\optsv}[4]{\par\vspace{3pt}\begin{tabular}{@{}l@{\ \ }p{0.9\linewidth}@{}}(A) & #1\\(B) & #2\\(C) & #3\\(D) & #4\end{tabular}\par\vspace{9pt}}

\fancyhead[R]{\footnotesize\color{black!60}Unit 1 test}
\begin{document}

\wstitle{1.11}{Unit 1 Test: Limits and Continuity}{Unit 1: Limits and Continuity \textperiodcentered{} 20 questions \textperiodcentered{} 45 minutes}

\begin{scorebox}{Instructions}
\begin{itemize}[leftmargin=*,itemsep=2pt,topsep=2pt]
  \item \textbf{Part A} (questions 1--15): 32 minutes, \textbf{no calculator}.
  \item \textbf{Part B} (questions 16--20): 13 minutes, a graphing calculator is required.
  \item Each question has exactly one correct answer. Circle one letter. There is no penalty for a wrong answer, so answer every question.
  \item Questions 2, 3 and 14 refer to the graph of $q$ printed with question 2.
\end{itemize}
\end{scorebox}

\sech{Part A \textperiodcentered{} No calculator}
\begin{enumerate}[leftmargin=*,itemsep=0pt]
  \item $\displaystyle\lim_{x\to 2}\dfrac{x^2-4}{x^2+x-6}=$
  \opts{$\dfrac45$}{$0$}{$1$}{does not exist}

  \item The graph of the function $q$, defined on $[-4,5]$, is shown.
\begin{apgraph}{-5}{6}{-1.5}{5.5}
\addplot[domain=-4:-2]{-x-2};
\addplot[domain=-2:1]{x+2};
\addplot[domain=1:5]{5-x};
\closeddot{-4}{2}\closeddot{-2}{2}\closeddot{1}{3}\closeddot{5}{0}
\opendot{-2}{0}\opendot{1}{4}
\node[font=\small, color=ink] at (axis cs:3.9,2.2) {$y=q(x)$};
\end{apgraph}
  What is $\displaystyle\lim_{x\to -2}q(x)$?
  \opts{$2$}{$0$}{$-2$}{does not exist}

  \item For the function $q$ graphed above, $\displaystyle\lim_{x\to 1}q(x)$ is
  \opts{$3$}{$4$}{$\dfrac72$}{nonexistent}

\newpage
  \item $\displaystyle\lim_{x\to 0}\dfrac{\sin(3x)}{4x}=$
  \opts{$0$}{$\dfrac34$}{$\dfrac43$}{$1$}

  \item $\displaystyle\lim_{x\to \infty}\dfrac{6x^2-1}{2x^2+5x}=$
  \opts{$0$}{$3$}{$6$}{$\infty$}

  \item $\displaystyle\lim_{x\to -\infty}\dfrac{\sqrt{16x^2+3}}{2x}=$
  \opts{$-2$}{$2$}{$8$}{$-8$}

  \item For what value of $k$ is $f(x)=\begin{cases}x^2+k,&x\le1\\3x+2k,&x>1\end{cases}$ continuous at $x=1$?
  \opts{$2$}{$0$}{$-2$}{$-1$}

  \item At $x=5$, the function $f(x)=\dfrac{x^2-25}{x^2-4x-5}$ has
  \optsv{a removable discontinuity}{a jump discontinuity}{an infinite discontinuity}{no discontinuity}

  \item The function $f$ is continuous on $[0,8]$. Selected values are shown.
\begin{center}\renewcommand{\arraystretch}{1.25}\begin{tabular}{|c|c|c|c|c|}\hline $x$ & $0$ & $2$ & $5$ & $8$ \\ \hline $f(x)$ & $3$ & $-4$ & $1$ & $-2$ \\ \hline\end{tabular}\end{center}
  What is the least number of zeros of $f$ in $[0,8]$?
  \opts{$1$}{$2$}{$3$}{$4$}

  \item $\displaystyle\lim_{x\to 9}\dfrac{\sqrt x-3}{x-9}=$
  \opts{$\dfrac16$}{$\dfrac13$}{$0$}{does not exist}

  \item If $\displaystyle\lim_{x\to 4}f(x)=2$ and $\displaystyle\lim_{x\to 4}g(x)=-1$, then $\displaystyle\lim_{x\to 4}\dfrac{3f(x)-[g(x)]^2}{f(x)+g(x)}=$
  \opts{$-5$}{$7$}{$\dfrac53$}{$5$}

\newpage
  \item A function $f$ satisfies $4-x^2\le f(x)\le4+x^2$ for all $x$. Then $\displaystyle\lim_{x\to 0}f(x)$
  \optsv{is $0$}{is $4$}{is $8$}{cannot be determined from this information}

  \item $\displaystyle\lim_{x\to 3^-}\dfrac{x+1}{x-3}=$
  \opts{$4$}{$\infty$}{$-\infty$}{$0$}

  \item For the function $q$ graphed in question 2, $\displaystyle\lim_{x\to 1^+}q\big(q(x)\big)=$
  \opts{$2$}{$3$}{$0$}{$1$}

  \item The graph of $f(x)=\dfrac{x-1}{x^2-1}$ has a vertical asymptote at
  \optsv{$x=-1$ only}{$x=1$ only}{$x=-1$ and $x=1$}{no value of $x$}
\end{enumerate}

\newpage
\sech{Part B \textperiodcentered{} Graphing calculator required}
\begin{enumerate}[leftmargin=*,itemsep=0pt,start=16]
  \item To three decimal places, $\displaystyle\lim_{x\to 0}\dfrac{e^{2x}-1}{\sin x}$ is
  \opts{$1.000$}{$0.500$}{$2.000$}{$0.035$}

  \item To three decimal places, $\displaystyle\lim_{x\to 0}\dfrac{1-\cos(3x)}{x^2}$ is
  \opts{$9.000$}{$4.500$}{$1.500$}{$0.000$}

  \item Selected values of a function $g$ are shown.
\begin{center}\renewcommand{\arraystretch}{1.25}\begin{tabular}{|c|c|c|c|c|c|}\hline $x$ & $1.9$ & $1.99$ & $2$ & $2.01$ & $2.1$ \\ \hline $g(x)$ & $4.61$ & $4.9601$ & $1$ & $5.0401$ & $5.41$ \\ \hline\end{tabular}\end{center}
  The best estimate of $\displaystyle\lim_{x\to 2}g(x)$ is
  \opts{$1$}{$4.96$}{nonexistent}{$5$}

  \item To three decimal places, $\displaystyle\lim_{x\to 1}\dfrac{\ln x}{x-1}$ is
  \opts{$1.000$}{$0.000$}{$2.718$}{$0.368$}

  \item A function $h$ is continuous on $[1,4]$ with $h(1)=-2.3$ and $h(4)=1.7$. Which statement must be true?
  \optsv{$h$ is increasing on $[1,4]$.}{$h(2.5)=-0.3$.}{$h(c)=0$ for some $c$ in $(1,4)$.}{$h$ has exactly one zero in $(1,4)$.}
\end{enumerate}

\end{document}
